problem:  
Let \((M_i^n, g_i)\) be a non-collapsing sequence of closed Riemannian manifolds with 
\(\operatorname{Ric}(M_i) \geq -(n-1)\) and \(\operatorname{diam}(M_i) \le D\).  
Assume that \((M_i, g_i)\) converge in the measured Gromov–Hausdorff sense to a compact limit space \((X, d_X, \mathcal{H}^n)\).  
Cheeger–Colding theory implies the existence of a regular set \(\mathcal{R} \subset X\), which is an open dense subset that is \(C^{1,\alpha}\)-rectifiable and carries a smooth Riemannian structure.  

The conjectural question is:  
Does there exist (after possibly passing to a subsequence) a sequence of maps  
\[
f_i : M_i \to X
\]  
realizing the Gromov–Hausdorff convergence such that:  
1. For any compact \(K \subset \mathcal{R}\), there exist open subsets \(U_i \subset M_i\) with \(f_i(U_i)\supset K\) on which \(f_i\) is an almost Riemannian submersion in the sense that  
   \[
   (1-\varepsilon_i)\, |v|_{g_i} \le |df_i(v)|_{d_X} \le (1+\varepsilon_i)\, |v|_{g_i} \quad \text{for all } v\in T(U_i),
   \]
   with \(\varepsilon_i \to 0\);  
2. The complement \(M_i \setminus U_i\) has volume tending to zero as \(i\to\infty\).  

Equivalently, can the smooth approximating manifolds \(M_i\) be globally decomposed—up to small volume error—as *almost Riemannian coverings* of the regular set of the limit \(X\)?  
If true, such a statement would provide a geometric-topological lifting of Colding’s volume convergence theorem to a fibering-type convergence respecting the infinitesimal Riemannian structures.

Why is it a "good" problem:  
This reformulated problem is mathematically well-posed: it avoids undefined notions like “branched coverings of metric spaces” and instead uses the standard framework of measured Gromov–Hausdorff convergence and almost Riemannian submersions on the regular part of the limit. It directly probes a frontier question in the analysis of Ricci limit spaces—whether one can *upgrade metric convergence to a differentiable fibering structure* on large-measure subsets. This engages several deep themes: the structure and topology of the singular set, the stability of fundamental groups under Ricci bounds, and the local Riemannian nature of limits. If resolved, such a result would form a bridge between geometric analysis, metric geometry, and geometric topology, refining Colding’s and Cheeger–Colding’s foundational convergence theorems and pushing them toward a new structural classification of Ricci limit spaces. The statement is compact, conceptually clear, but profoundly deep—exemplifying the “narrow entrance, profound depth” ideal of a good mathematical problem.